% 原书 p.86–p.88（PDF p.94–p.96）
% 第 4 章 4.1 节：Moser 技巧

\chapter{Moser 论证（Moser's Argument）}

% >>> p094
Moser 论证（Moser's argument，亦称 Moser 技巧）是一种揭示辛结构与切触结构「柔性」的方法。
这种柔性常常有助于把问题化归到标准模型上讨论。这里的标准模型包括 Euclid 空间
（如例 1.2.1 与例 3.2.1），以及余切丛（如例 1.2.2）及其切触化（contactization，如例 3.2.2）。
在处理局部信息时，这样做可以大幅简化讨论或计算。本章将展示 Moser 论证及其各种应用。

\section{Moser 技巧（Moser's Trick）}

\begin{keytakeaways}
\begin{itemize}
  \item 几何图像：若一族辛形式沿参数方向的变化都可由恰当 $2$-形式实现，就能用一族辛同胚把它
        拉回到初始形式；切触情形只需给定 $\xi_t=\ker\alpha_t$，条件更宽。
  \item 数学工具：Cartan 公式与流的求导、Reeb 向量场分解
        $X_t=a(t)R_{\alpha_t}+Y_t$（$Y_t\in\xi_t$）。
  \item 本节结论：Moser 技巧（定理 4.1.1）与 Gray 定理（定理 4.1.2）；但切触 $1$-形式一般
        不能沿同痕保持（练习 4.1.1）。
\end{itemize}
\end{keytakeaways}

Moser 技巧源自 Moser 的论文 [164]，其目标是通过一族辛同胚把两个辛结构联系起来。
具体地说，设 $M$ 是紧辛流形，$\{\omega_t\}_{t\in[0,1]}$ 是 $M$ 上的一族光滑辛结构，
并假设存在一族光滑 $1$-形式 $\{\sigma_t\}_{t\in[0,1]}$，使得
\begin{equation}\tag{4.1}
  \frac{\dd\omega_t}{\dd t}=\dd\sigma_t.
\end{equation}
那么就能找到 $M$ 上一族微分同胚的同痕 $\{\varphi_t\}_{t\in[0,1]}$，使得
$\varphi_t^*\omega_t=\omega_0$。特别地，$(M,\omega_0)$ 与 $(M,\omega_1)$ 经由辛同胚
$\varphi_1$ 而辛同构。

一如往常，要得到这样的同痕 $\{\varphi_t\}_{t\in[0,1]}$，我们首先需要找到一个向量场 $X_t$
（它积分后给出该同痕，即流 $\{\varphi_t\}_{t\in[0,1]}$）。$X_t$ 应满足的条件可由一个
「反向」论证得到，而这个论证最终来自条件 (4.1)。假设这样的 $\{\varphi_t\}_{t\in[0,1]}$ 存在，
则
\begin{equation*}
  \begin{aligned}
    0&=\frac{\dd}{\dd t}\bigl(\varphi_t^*\omega_t\bigr)
    =\varphi_t^*\Bigl(\frac{\dd\omega_t}{\dd t}\Bigr)
    +\varphi_t^*\bigl(\mathcal L_{X_t}\omega_t\bigr)\\
    &=\varphi_t^*\bigl(\dd\sigma_t+\dd\iota_{X_t}\omega_t\bigr)\\
    &=\dd\bigl(\varphi_t^*(\sigma_t+\iota_{X_t}\omega_t)\bigr).
  \end{aligned}
\end{equation*}
% [Tutor Note: 易错预警——第一行到第二行用了 Cartan 公式
% L_{X_t}ω_t = d ι_{X_t}ω_t + ι_{X_t} d ω_t，而每个 ω_t 都是闭的（辛形式），故后一项消失；
% 第二行到第三行用了 φ_t^* 与 d 可交换（φ_t 是微分同胚）。另请注意 X_t 是「同痕的生成向量场」，
% 约定 dφ_t/dt = X_t∘φ_t；若把 X_t 换成 -X_t，得到的流是逆同痕，核对时应以 φ_t^*ω_t = ω_0
% 的归一化条件为准。]
例如，当 $\sigma_t+\iota_{X_t}\omega_t=0$ 时上式成立。一旦 (4.1) 中的 $\sigma_t$ 给定，
由 $\omega_t$ 的非退化性，向量场 $X_t$ 便由方程 $\sigma_t+\iota_{X_t}\omega_t=0$ 唯一确定。

% >>> p095
注意我们需要 $\sigma_t$ 具有一定的正则性（光滑性），使所得的向量场 $X_t$ 足够光滑、可以积分
（参见定理 4.2.2 的证明）。总之，我们刚刚证明了下面这个著名结果。

\begin{wthm}{4.1.1}（Section 3.2 in [153]）
  设 $(M,\omega_0)$ 是紧辛流形，$\{\omega_t\}_{t\in[0,1]}$ 是从 $\omega_0$ 出发的一族光滑
  辛结构，并设存在一族光滑 $1$-形式 $\{\sigma_t\}_{t\in[0,1]}$ 使得
  $\dd\omega_t/\dd t=\dd\sigma_t$。则存在辛同胚的同痕 $\{\varphi_t\}_{t\in[0,1]}$，
  使得 $\varphi_t^*\omega_t=\omega_0$。
\end{wthm}
% [Tutor Note: 含糊点拨——定理要求 M 紧，实质上是保证向量场 X_t 能在 t∈[0,1] 上整体积分成
% 流（紧性给出流的存在时间的一致性）；σ_t 的光滑性则保证 X_t 光滑，从而流光滑依赖于初值。
% 若去掉紧性，结论须改为局部成立（参看 4.3 节 Lemma 4.3.1 的证明，那里是在 Q 的邻域内做）。]

\begin{wrem}{4.1.1}
  同样的论证对任意闭定向流形上体积形式的类似形变 $\{\sigma_t\}_{t\in[0,1]}$ 也成立，
  其中 $\sigma_t$ 既非退化又闭。
\end{wrem}
% [Tutor Note: 原文疑误（记号/条件含混）——此处以 σ_t 记「体积形式族」，与 (4.1) 中的
% 1-形式 σ_t 记号冲突；且「非退化且闭」是辛形式的刻画，并非一般体积形式的性质。结合
% 定理 4.2.1 的准确表述，作者的意思应是：把两条辛形式换成两个体积形式、把沿 t 的导数写成
% 恰当形式，论证照样成立。译文按原文照译，详见 qsg/notes/review_ch04.md。]

在切触几何中，与 (4.1) 类似的条件甚至更简单。事实上，对切触结构几乎没有附加条件。
设 $M$ 是紧切触流形，$\{\xi_t\}_{t\in[0,1]}$ 是 $M$ 上的一族切触结构，并假设存在一族光滑
切触 $1$-形式 $\{\alpha_t\}_{t\in[0,1]}$，使得对任意 $t\in[0,1]$ 都有
$\xi_t=\ker\alpha_t$。与上面类似，可以找到一族微分同胚的同痕 $\{\psi_t\}_{t\in[0,1]}$，
它们保持切触结构（但不一定保持切触 $1$-形式！），使得 $(\psi_0)_*\xi_0=\xi_t$。
特别地，$(M,\xi_0)$ 与 $(M,\xi_1)$ 经由切触同胚 $\psi_1$ 彼此切触同构。
% [Tutor Note: 原文疑误——此处应为 (ψ_t)_*ξ_0 = ξ_t。因为 ψ_0 = id，左端恒为 ξ_0；
% 除非 ξ_t 不随 t 变化，否则 (ψ_0)_*ξ_0 = ξ_t 不可能对一般的 t 成立。译文保留原文写法，
% 详见 qsg/notes/review_ch04.md。]

在形式层面，所需的同痕 $\{\psi_t\}_{t\in[0,1]}$ 满足
$\psi_t^*\alpha_t=\ee^{f_t}\alpha_0$，其中 $f_t\colon M\to\R$ 是一族光滑共形因子。
沿用辛情形的同样思路，我们要找某个向量场 $X_t$（它将积分成 $\{\psi_t\}_{t\in[0,1]}$）。
对 $\ee^{f_t}\alpha_0=\psi_t^*\alpha_t$ 两边求导，得
\begin{equation}\tag{4.2}
  \begin{aligned}
    \ee^{f_t}\cdot\frac{\dd f_t}{\dd t}\alpha_0
    &=\psi_t^*\frac{\dd\alpha_t}{\dd t}+\psi_t^*\bigl(\mathcal L_{X_t}\alpha_t\bigr)\\
    &=\psi_t^*\frac{\dd\alpha_t}{\dd t}
    +\psi_t^*\bigl(\dd(\alpha_t(X_t))+\iota_{X_t}\dd\alpha_t\bigr).
  \end{aligned}
\end{equation}
把 $\alpha_0$ 换成 $\ee^{-f_t}\psi_t^*\alpha_t$，则 (4.2) 可改写为
\begin{equation*}
  \ee^{f_t}\cdot\frac{\dd f_t}{\dd t}\cdot\bigl(\ee^{-f_t}\psi_t^*\alpha_t\bigr)
  =\psi_t^*\frac{\dd\alpha_t}{\dd t}
  +\psi_t^*\bigl(\dd(\alpha_t(X_t))+\iota_{X_t}\dd\alpha_t\bigr),
\end{equation*}
亦即
\begin{equation*}
  \frac{\dd f_t}{\dd t}\psi_t^*\alpha_t
  =\psi_t^*\Bigl(\frac{\dd\alpha_t}{\dd t}
  +\dd(\alpha_t(X_t))+\iota_{X_t}\dd\alpha_t\Bigr).
\end{equation*}
于是需要求解一个向量场 $X_t$，使它满足如下方程：
\begin{equation}\tag{4.3}
  \Bigl(\frac{\dd f_t}{\dd t}\circ\psi_t^{-1}\Bigr)\cdot\alpha_t
  =\frac{\dd\alpha_t}{\dd t}+\dd(\alpha_t(X_t))+\iota_{X_t}\dd\alpha_t.
\end{equation}
仿照图 3.1 中的分解，写 $X_t=a(t)R_{\alpha_t}+Y_t$，其中 $a(t)$ 是某个光滑函数，
$Y_t\in\xi_t$ 是某个向量场（待定），则
\begin{equation*}
  \dd(\alpha_t(X_t))=\dd a(t)
  \quad\text{和}\quad
  \iota_{X_t}\dd\alpha_t=\dd\alpha_t(Y_t,\cdot).
\end{equation*}
% [Tutor Note: 易错预警——d(α_t(X_t)) = da(t) 中的 a(t) 是 M 上的光滑函数（t 只是参数，
% 应读作 a(t,·)），不是关于 t 的一元函数；da(t) 是对空间变量取外微分。又 R_{α_t} 是 α_t 的
% Reeb 向量场，满足 α_t(R_{α_t}) = 1、ι_{R_{α_t}}dα_t = 0，而 Y_t ∈ ξ_t = ker α_t；
% 这正是下文 (4.4)—(4.6) 能逐项化简的依据。]

% >>> p096
因此，(4.3) 简化为
\begin{equation}\tag{4.4}
  \Bigl(\frac{\dd f_t}{\dd t}\circ\psi_t^{-1}\Bigr)\cdot\alpha_t
  =\frac{\dd\alpha_t}{\dd t}+\dd a(t)+\dd\alpha_t(Y_t,\cdot).
\end{equation}
然后把 $R_{\alpha_t}$ 代入 (4.4)，得到
\begin{equation}\tag{4.5}
  \frac{\dd f_t}{\dd t}\circ\psi_t^{-1}
  =\frac{\dd\alpha_t}{\dd t}(R_{\alpha_t})+\dd a(t)(R_{\alpha_t}).
\end{equation}
(4.5) 中似乎有两个未知量 $f_t$ 与 $a(t)$，但请记住：为了得到所求的 $\psi_t$，我们只需要
找到一个 $X_t$。因此，若取 $a(t)\equiv0$，从而 $X_t=Y_t\in\xi_t$，便可由
\begin{equation*}
  \frac{\dd f_t}{\dd t}\circ\psi_t^{-1}
  =\frac{\dd\alpha_t}{\dd t}(R_{\alpha_t})+0
\end{equation*}
算出 $f_t$。此外，(4.4) 简化为
\begin{equation}\tag{4.6}
  \frac{\dd\alpha_t}{\dd t}(R_{\alpha_t})\cdot\alpha_t
  -\frac{\dd\alpha_t}{\dd t}=\dd\alpha_t(Y_t,\cdot),
\end{equation}
由 $\dd\alpha_t|_{\xi_t}$ 的非退化性（来自 $1$-形式
$\frac{\dd\alpha_t}{\dd t}(R_{\alpha_t})\cdot\alpha_t-\frac{\dd\alpha_t}{\dd t}$
在 $\xi_t$ 上的限制），上式可解出 $Y_t$（从而解出向量场 $X_t$）。
% [Tutor Note: 易错预警——(4.6) 左端与右端都是 1-形式；在 ξ_t 上取 Z ∈ ξ_t 即化为
% 标量方程 (dα_t/dt)(R_{α_t})·α_t(Z) − (dα_t/dt)(Z) = dα_t(Y_t,Z)。注意 α_t(Z) = 0，
% 故第一项在 ξ_t 上为零，于是 dα_t(Y_t,Z) = −(dα_t/dt)(Z)；由 dα_t|_{ξ_t} 非退化
% （ξ_t 是 2n 维切触分布）唯一确定 Y_t ∈ ξ_t。]

\begin{wrem}{4.1.2}
  当然，也可以取另一个 $a(t)$，比如取 $a(t)\equiv c$，其中 $c\neq0$ 为常数。此时 (4.6)
  仍然成立，得到的 $Y_t$ 与上面 $a(t)\equiv0$ 的情形相同。但此时
  $X_t=cR_{\alpha_t}+Y_t$，它积分得的同痕 $\{\psi_t\}_{t\in[0,1]}$ 可能具有不同的共形因子。
\end{wrem}

总之，我们证明了另一个著名结果。

\begin{wthm}{4.1.2}（Gray [98]，亦见 [92] 中的定理 2.2.2）
  设 $(M,\xi_0)$ 是紧切触流形，$\{\xi_t\}_{t\in[0,1]}$ 是从 $\xi_0$ 出发的一族切触结构，
  且存在一族光滑切触 $1$-形式 $\{\alpha_t\}_{t\in[0,1]}$ 使得 $\xi_t=\ker\alpha_t$。则存在
  切触同胚的光滑同痕 $\{\psi_t\}_{t\in[0,1]}$，使得 $(\psi_t)_*\xi_0=\xi_t$。
\end{wthm}
% [Tutor Note: 术语/背景——Gray 定理只保证「切触结构」沿同痕不变，而不保证切触 1-形式不变：
% 保持切触结构的同痕只把 α_t 拉回到 e^{f_t}α_0（共形因子可依赖 t 与空间点）。若要求
% ψ_t^*α_t = α_0（严格切触同胚意义下的同痕），一般做不到，这正是练习 4.1.1 的要点。
% 定理要求 M 紧，原因与定理 4.1.1 相同。]

\begin{wexer}{4.1.1}
  举一个例子说明 Gray 定理对切触 $1$-形式不成立。具体地说，构造一族光滑切触 $1$-形式
  $\{\alpha_t\}_{t\in[0,1]}$，使得不存在光滑同痕 $\{\varphi_t\}_{t\in[0,1]}$ 满足
  $\varphi_t^*\alpha_t=\alpha_0$。
\end{wexer}
